\section{Introduction}\label{sec:intro}
Software engineering is rapidly delegating decisions to automated tools: which code needs a security review, which alerts can be ignored, which patch can be merged. Each delegation shifts a question from ``is the tool effective?'' to ``can its output be trusted without inspection?''. For vulnerability detection the answer matters twice: a false alarm wastes reviewer time, but a missed vulnerability is exactly the event that the automation was supposed to prevent, and it is invisible by construction when the tool clears the code.

The research literature on learning-based vulnerability detection reports accuracy, F1 or AUROC, mostly on random splits of public corpora. Such numbers do not answer the question a security team has to answer before it stops reviewing the cleared code: \emph{what fraction of the vulnerable functions will the tool let through, and does that fraction remain valid on the next project, the next release, or another codebase?} Recent studies have shown that accuracy collapses on realistic data and that the corpora themselves are noisy and duplicated~\cite{chakraborty2021deep,croft2023data,ding2024primevul}, but they stop at measuring detection quality. Two parts of the problem remain open: (i) a way to \emph{control} the miss rate of the decisions that are actually delegated, and (ii) an empirical account of whether that control survives the distribution shifts that characterise software evolution.

\paragraph{Approach} We treat clearing a function as a decision that must carry a statistical guarantee, and build it on conformal risk control~\cite{vovk2005algorithmic,angelopoulos2023conformal}. The detector is used only as a scorer; a threshold chosen on labelled calibration positives bounds the share of vulnerable functions that are auto-cleared by a user-chosen budget $\alpha$ (Figure~\ref{fig:framework}). The rest of the code goes to human review, or is auto-flagged where a precision guarantee can be certified. We complement the construction with the quantities needed to deploy it: a high-probability version, group-conditional calibration, a bound on how far distribution shift can push the realised miss rate, a cost model for choosing $\alpha$, and a monitoring/re-calibration loop.

\begin{figure}[t]
\centering
\resizebox{\linewidth}{!}{\begin{tikzpicture}[font=\footnotesize,
  box/.style={draw, rounded corners=2pt, align=center, inner sep=3pt, fill=#1, minimum height=9mm},
  box/.default=white, arr/.style={-{Latex[length=2mm]}, thick}]
  \node[box=blue!8] (fn) at (0,0) {Function\\source $x$};
  \node[box=blue!8] (det) at (2.6,0) {Detector\\score $s(x)$};
  \node[box=orange!15, text width=33mm] (p) at (6.6,0) {Conformal $p$-value\\[1pt]$p(x)=\dfrac{1+\#\{i:\,s_i\le s(x)\}}{n+1}$};
  \node[box=green!14, text width=27mm] (clear) at (11.4,1.5) {\textbf{Auto-clear}\\$p(x)\le\alpha$\\miss rate $\le\alpha$};
  \node[box=gray!14, text width=27mm] (rev) at (11.4,0) {\textbf{Human review}\\otherwise};
  \node[box=red!10, text width=27mm] (fl) at (11.4,-1.5) {\textbf{Auto-flag}\\$s(x)\ge\hat\lambda$\\precision certified};
  \node[box=white, text width=40mm] (cal) at (6.6,-2.9) {Calibration positives $\{s_i\}_{i\le n}$\\from the \emph{deployment} distribution};
  \node[box=yellow!15, text width=36mm] (mon) at (1.4,-2.9) {Monitor outcomes;\\re-calibrate with $k$ new labelled positives};
  \draw[arr] (fn)--(det); \draw[arr] (det)--(p);
  \draw[arr] (p.east) -- ++(0.5,0) |- (clear.west);
  \draw[arr] (p.east) -- (rev.west);
  \draw[arr] (p.east) -- ++(0.5,0) |- (fl.west);
  \draw[arr] (cal)--(p);
  \draw[arr, dashed] (mon)--(cal);
  \draw[arr, dashed] (clear.east) -- ++(0.55,0) -- ++(0,-5.3) -- ++(-12.2,0) -- (mon.south);
\end{tikzpicture}
}
\caption{Risk-controlled triage. Calibration positives from the deployment distribution fix the clearance threshold; a function is auto-cleared when its conformal $p$-value is at most the miss budget $\alpha$, auto-flagged when a certified precision threshold is met, and reviewed otherwise. Outcomes feed back into re-calibration.}\label{fig:framework}
\end{figure}

\paragraph{Findings in brief} We test the construction on 487\,412 deduplicated C/C++ functions from two public corpora (DiverseVul and Big-Vul), in seven evaluation scenarios (random, project-disjoint, temporal, cross-dataset), with three seeds and four detectors, together with a rule-based analyser. The guarantee is exact on exchangeable data and fails under every kind of shift we examine, by a factor of 2.7 on average; the standard remedy of covariate-shift weighting does not repair it; and a small number of labelled target positives (about 100) does. What can then be automated is about half of what a random-split evaluation suggests.

\paragraph{Contributions}
\begin{enumerate}[itemsep=1pt]
\item \textbf{Formulation and theory.} A risk-control formulation of vulnerability triage with operational quantities (miss rate, safe automation rate, review effort) and the following results: validity and calibration-set variability of conformal clearance (Prop.~\ref{prop:valid}, Cor.~\ref{cor:beta}), a high-probability variant (Prop.~\ref{prop:pac}), a bound on the miss-rate excess under shift in terms of the KS distance (Prop.~\ref{prop:shift}), group-conditional validity (Prop.~\ref{prop:mondrian}), the identity $\mathrm{AUROC}=\int\mathrm{TNR}(\alpha)d\alpha$ that explains why AUROC is a poor proxy for operational value (Lemma~\ref{lem:auc}), and an economically optimal miss budget (Prop.~\ref{prop:cost}).
\item \textbf{A shift-aware empirical test of the guarantee} on seven scenarios, including evidence that source-calibrated clearance misses $2.7\times$ the nominal share under shift, that all 336 observations satisfy the proposed bound, that covariate-shift weighting is insufficient, and how many labelled target positives are needed to recover validity.
\item \textbf{A data-quality and leakage audit} of DiverseVul and Big-Vul (24.6\% exact duplicates in Big-Vul; 42.7\% of Big-Vul inside DiverseVul) and its quantified effect on cross-dataset evaluation.
\item \textbf{Accountability analyses}: uneven miss rates across function sizes and CWEs and their repair by Mondrian calibration; decision flips under semantics-preserving rewrites; a TreeSHAP fidelity/consistency study that links the model's reliance on function size to the unevenness.
\item \textbf{Operational evidence}: comparison with a rule-based analyser and a cost analysis that shows when delegation pays off.
\item \textbf{An open, reproducible artifact} that regenerates every number in this paper from cached scores.
\end{enumerate}

\paragraph{Scope} Because the study had to be completed on CPU-only infrastructure, the detectors are TF-IDF/gradient-boosting models rather than fine-tuned code language models. The triage layer is detector-agnostic and the theoretical conclusions do not depend on the scorer, but the numerical levels (for example the achievable SAR) will differ for stronger models; Section~\ref{sec:threats} discusses this.

\paragraph{Structure} Section~\ref{sec:related} places the work in the literature, Section~\ref{sec:method} develops the theory, Section~\ref{sec:design} describes the study, Section~\ref{sec:results} answers the five research questions, and Sections~\ref{sec:discussion}--\ref{sec:conclusion} discuss implications, threats and conclusions.
